\uniterablesection{Вступ}

Звіт відображає результати передкваліфікаційної практики, метою якої було
опрацювання та проєктування напрямів удосконалення інформаційної системи моніторингу
особистих звичок Dailu. Базова версія системи частково знімає навантаження ручного обліку за
рахунок автоматичного збору активності із зовнішніх сервісів GitHub і Strava, проте зібрані
дані використовуються пасивно, а активність поза межами цифрових сервісів не охоплюється.

Під час практики опрацьовано два напрями розвитку системи, які й становлять зміст цього звіту.
Перший — інтелектуальна обробка накопичених даних засобами штучного інтелекту: генерування
персоналізованих рекомендацій та інсайтів, автоматична категоризація звичок і діалогова
взаємодія з користувачем через мовну модель. Другий — розширення джерел активності у фізичний
світ за допомогою простого пристрою Інтернету речей (IoT), що дає змогу фіксувати виконання
звичок безпосередньою фізичною дією, а не лише через програмні події.

\textbf{Актуальність теми.} Ринок застосунків для трекінгу звичок стрімко зростає й переходить
до інтелектуальних, адаптивних рішень: сучасні застосунки дедалі частіше інтегрують засоби
штучного інтелекту для персоналізації досвіду користувача. Паралельно у сфері цифрового
здоров'я та зміни поведінки зростає інтерес до застосування великих мовних моделей для
персоналізованого коучингу, а поширення доступних сенсорних і носимих пристроїв Інтернету речей
розширює можливості об'єктивного збору даних про поведінку людини \citecustom{3, 4}. Тема
роботи знаходиться саме на перетині цих тенденцій.

\textbf{Мета практики} — опрацювати предметну область та спроєктувати напрями вдосконалення
системи Dailu шляхом інтеграції технологій штучного інтелекту для інтелектуального аналізу
поведінки користувача та пристрою Інтернету речей для автоматизованого збору активності з
фізичного світу.

Для досягнення поставленої мети під час практики розв'язано такі завдання:

\begin{enumerate}[noitemsep, topsep=2pt, leftmargin=*, label=\arabic*)]
    \item проаналізувати предметну область інтелектуального моніторингу звичок, дослідити наявні рішення-аналоги та методи застосування штучного інтелекту й Інтернету речей у зміні поведінки;
    \item виявити обмеження наявних рішень і обґрунтувати доцільність удосконалення;
    \item сформулювати функціональні та нефункціональні вимоги до вдосконаленої системи;
    \item обґрунтувати архітектуру AI-підсистеми (рекомендації, автокатегоризація, діалоговий асистент) з урахуванням наявної модульної архітектури;
    \item обґрунтувати апаратно-програмну взаємодію IoT-пристрою із серверною частиною системи;
    \item уточнити технологічний стек удосконаленої системи.
\end{enumerate}

\textbf{Об'єкт дослідження} — процес формування, автоматизованого обліку та інтелектуального
аналізу особистих звичок користувача.

\textbf{Предмет дослідження} — методи та програмні засоби застосування штучного інтелекту й
технологій Інтернету речей для збору, обробки та інтелектуального аналізу даних про активність
користувача.

\textbf{Методи.} Застосовано методи системного аналізу — для вивчення предметної області та
декомпозиції задачі; принципи предметно-орієнтованого проєктування й чистої архітектури — для
побудови та розширення доменної моделі; методи машинного навчання й обробки природної мови — як
основу інтелектуальних функцій системи.

\textbf{Структура звіту.} Звіт містить календарний план проходження практики, вступ, два
розділи, висновки та список використаних джерел. У першому розділі проаналізовано предметну
область, здійснено огляд літератури й наявних рішень та сформульовано вимоги до вдосконаленої
системи. Другий розділ присвячено обґрунтуванню та проєктуванню нових складових — AI-підсистеми
та IoT-каналу збору активності.
